\section{Conclusion}
\label{sec:conclusion}

We reported a controlled descriptive study of authentic ISV corpus priming for Polish-to-ISV LLM translation across HIGH, LOW, and UNSEEN source regimes.
Matched-six mean canonical-coverage $\Delta$ was $+7.39$\,pp (HIGH), $+8.53$\,pp (LOW), and $+5.78$\,pp (UNSEEN).
Effects were heterogeneous across configurations, including Gemini sign reversals on HIGH versus LOW/UNSEEN.
Qualitative audits linked some HIGH positives to corpus-opening reuse, while LOW/UNSEEN positives often showed lexical and orthographic shifts without winter-narrative copying.
One Qwen LOW pairing was invalid and excluded.
Within these limits---$n{=}3$, descriptive analysis, coverage rather than full quality---the evidence supports a restrained claim: corpus priming is associated with generally positive but configuration-dependent coverage shifts that can persist beyond high-overlap narrative sources.


Future revisions may add human evaluation, additional unseen-domain sources with orthogonalized length, and verified Interslavic language-description citations.
Those extensions would strengthen quality and linguistic framing; they are not required to interpret the descriptive coverage contrasts already measured here.
